\documentclass{article}

% if you need to pass options to natbib, use, e.g.:
%     \PassOptionsToPackage{numbers, compress}{natbib}
% before loading neurips_2025

% ready for submission
\usepackage{neurips_2025}

% to compile a preprint version, e.g., for submission to arXiv, add add the
% [preprint] option:
%     \usepackage[preprint]{neurips_2025}

% to compile a camera-ready version, add the [final] option, e.g.:
%     \usepackage[final]{neurips_2025}

% to avoid loading the natbib package, add option nonatbib:
%    \usepackage[nonatbib]{neurips_2025}

% Packages required by the reference template
\usepackage[utf8]{inputenc} % allow utf-8 input
\usepackage[T1]{fontenc}    % use 8-bit T1 fonts
\usepackage{hyperref}       % hyperlinks
\usepackage{url}            % simple URL typesetting
\usepackage{booktabs}       % professional-quality tables
\usepackage{amsfonts}       % blackboard math symbols
\usepackage{nicefrac}       % compact symbols for 1/2, etc.
\usepackage{microtype}      % microtypography
\usepackage{xcolor}         % colors

% Additional packages used in the paper
\usepackage{graphicx}
\usepackage{caption}
\usepackage{subcaption}
\usepackage{amsmath}
\usepackage{amssymb}
\usepackage{bm}

\title{MoE-Harmony: Scalable Mixture-of-Experts Harmonization for Targeted, Interpretable Batch Correction in Single-Cell Transcriptomics}

\author{%
  David S.~Hippocampus\\
  Department of Computer Science\\
  Cranberry-Lemon University\\
  Pittsburgh, PA 15213 \\
  \texttt{hippo@cs.cranberry-lemon.edu} \\
}

\begin{document}

\maketitle

\begin{abstract}
Integrating large, heterogeneous single‑cell transcriptomic collections requires removing annotated technical covariates without erasing orthogonal biological structure—a task made difficult by sparse, heteroskedastic, and often nonlinear distortions and by frequent confounding between batch and biology. We propose Mixture‑of‑Experts Harmony (MoE‑Harmony), a principled, scalable harmonization algorithm that couples cosine‑normalized soft clustering with cluster‑wise ridge regression to produce state‑specific, targeted corrections. MoE‑Harmony preserves angular geometry via unit‑norm centroids, returns per‑expert regression coefficients for interpretability, and admits closed‑form, matrix‑level centroid, responsibility and ridge updates that support randomized block‑wise execution and memory‑efficient operation on commodity hardware. Across multi‑chemistry human PBMC collections and a large (~52k cell) mouse atlas we evaluate targeted correction fidelity (MTSCF), local mixing (LISI), responsibility entropy, covariance dissimilarity, and leave‑one‑donor robustness. Empirically, MoE‑Harmony concentrates correction energy in technical covariate subspaces, improves local batch mixing while reducing overcorrection in mixed neighborhoods, yields clearer expert specialization, and converges in fewer iterations than a strong baseline. An open reference implementation maps algorithmic steps to code to facilitate reproducibility and per‑expert audits.
\end{abstract}

\section{Introduction}

Integration of single‑cell transcriptomic datasets requires harmonizing either a high‑dimensional feature‑by‑cell matrix or a precomputed low‑dimensional embedding together with structured covariates that encode batch, donor, protocol and modality information for downstream joint analysis \citep{Liu2021RobustIO,Korsunsky2019FastSA}. Common preprocessing steps (centering, scaling, dimensionality reduction and columnwise normalization) set the stage for tasks such as cross‑study cell‑type discovery, marker inference, differential testing, trajectory analysis and spatial mapping.

The principal challenge is that technical variation (library chemistry, capture efficiency, sequencing depth and platform biases) produces sparse, heteroskedastic and often nonlinear distortions of expression space, while rich biological heterogeneity (donor effects, continuous trajectories, rare subpopulations, condition‑driven shifts) is frequently confounded with batch \citep{Svensson2017ExponentialSO,Hie2019EfficientIO,Grn2016DeNP,Richards2021ACO}. This confounding creates a dual failure mode for integration: over‑correction that erases true biology or under‑correction that preserves technical artefacts, motivating methods that can target annotated technical signal without collapsing orthogonal biological axes.

Existing algorithmic families embody contrasting separability assumptions and attendant failure modes. Regression‑based, empirical‑Bayes and design‑matrix corrections assume global, per‑feature additive or scaling effects and thus can erase cell‑state–specific biology when technical covariates are confounded with biology \citep{Li2022FastIntegrationAV,Lee2023BenchmarkingAF}. Neighbor‑matching and graph‑alignment approaches align local neighborhoods across batches but depend on adequate cross‑batch representation and preserved local geometry, conditions often violated when sampling densities or technical distortion vary \citep{Polaski2019BBKNNFB,Zhou2023AccurateIO}. Nonlinear manifold‑alignment and latent‑variable models can capture complex distortions but tend to produce corrections that are hard to interpret, unstable under confounding, and liable to distort covariance and gene–gene relationships \citep{McInnes2018UMAPUM,Xiong2021ConstructionOC,Stegle2012UsingPE,Xu2023MASIEF}.

Practical use further exposes three recurrent bottlenecks. First, scalability: dense linear algebra, neighbor graphs and per‑cluster solves inflate memory and runtime on datasets that now routinely reach millions of cells, forcing downsampling or aggregation strategies that sacrifice fidelity \citep{Bilous2021MetacellsUL,Hie2024ScanoramaIL}. Second, multimodal mismatches: dissociated scRNA‑seq, single‑nucleus and spatial assays differ in feature footprints and sampling biases so that a single homogeneous transform cannot reconcile modality‑specific covariance structure \citep{Azizi2018SingleCellMO,Galow2020IntegrativeCA,Baron2016AST}. Third, benchmarking gaps: many evaluations report either batch mixing or biological conservation in isolation rather than joint metrics that detect under‑ versus over‑correction, motivating standardized, covariate‑aware benchmarks \citep{Lein2007GenomewideAO,Der2019TubularCA,Rocque2021CreationOA}.

From these observations we distill practical desiderata: (i) targeted removal of annotated covariate‑associated signal while preserving orthogonal biological variation and gene–gene covariance; (ii) flexible, state‑specific models that allow covariate effects to vary across cell states; (iii) interpretable, per‑cluster diagnostics enabling audits of what was removed; and (iv) computational designs that exploit closed‑form, matrix‑local updates and blockwise streaming to operate on commodity hardware \citep{Zhang2019DefiningIC,Nicol2023RobustIO,Andreatta2020ProjectingST,Leek2007CapturingHI}. These desiderata emphasize conservative, state‑aware, diagnostically transparent and scalable harmonization.

A mixture‑of‑experts (MoE) approach naturally satisfies these requirements by coupling soft clustering with cluster‑wise regression: soft responsibilities accommodate continuous manifolds and reduce brittle hard assignments, cluster‑specific ridge solves permit state‑dependent covariate removal with explicit regularization and intercept control, and all core updates admit closed‑form, matrix‑based implementations that support randomized, block‑wise updates and numerical safeguards such as smoothing and clipping \citep{Dhillon2004ConceptDF,Zhang2023SynEcoSysAM}. Building on these principles, we design an MoE‑Harmony algorithm that localizes corrections to annotated covariate subspaces, preserves angular geometry via cosine normalization, returns per‑expert coefficients for interpretation, and implements memory‑efficient updates to scale to large single‑cell collections.

We evaluate this design on multi‑chemistry human PBMC collections and a large mouse atlas, projecting expression into PCA embeddings, encoding categorical covariates as design matrices and assessing harmonization with joint metrics that include subspace‑aligned correction fidelity, local mixing and covariance dissimilarity, as well as leave‑one‑donor robustness \citep{Zheng2016MassivelyPD,Arazi2019TheIC,PijuanSala2019ASM,Manno2018RNAVO,Segerstolpe2016SingleCellTP,Ritchie2015limmaPD}.

These considerations motivate the method and experiments that follow. Our primary contributions are summarized below:

\begin{itemize}
\item We propose Mixture‑of‑Experts Harmony (MoE‑Harmony), a principled harmonization algorithm that couples soft clustering of cosine‑normalized embeddings with cluster‑weighted ridge regression to produce cluster‑specific, targeted batch corrections while preserving angular geometry.
\item We introduce scalable algorithmic and implementation advances: matrix‑form centroid and responsibility updates, randomized block‑wise responsibility updates, closed‑form regularized ridge solves, and practical numerical safeguards to enable memory‑efficient operation on large single‑cell datasets.
\item We demonstrate empirically that MoE‑Harmony improves subspace‑aligned correction fidelity and interpretability across diverse benchmarks and metrics.
\item We release an open reference implementation mapping algorithmic steps to code to facilitate reproducibility and inspection of per‑expert corrections.
\end{itemize}

\section{Related Work}

\textbf{Single-cell transcriptomic data integration} Recent advances in this domain have employed probabilistic frameworks and deep learning models to effectively disentangle batch effects while preserving intrinsic biological heterogeneity. However, challenges persist in scalability, handling multiple covariates, and aligning multimodal data such as spatial transcriptomics. Existing methods like ComBat and mutual nearest neighbors (MNN) have demonstrated efficacy in mitigating batch-associated technical variation but often struggle to separate confounded biological heterogeneity when experimental covariates overlap \citep{Bttner2018ATM,Haghverdi2018BatchEI}. Embedding-based integration techniques employing adversarial learning or manifold alignment aim to project cells into shared latent spaces, addressing nonlinear batch effects but facing difficulties in interpretability and consistent preservation of cell-type-specific expression across heterogeneous datasets \citep{Liu2020JointlyDC,Regev2017TheHC}. Foundational work on defining cell types underscores the complexity posed by biological variability and technical confounders \citep{Regev2017TheHC}. Building upon this landscape, the proposed Mixture-of-Experts Harmony (MoE-Harmony) utilizes a mixture-of-experts architecture combined with efficient optimization strategies, enabling scalable integration across large datasets and multiple modalities while preserving nuanced biological variation~\citep{Maan2022TheDI,Dou2020UnbiasedIO}.

\textbf{Batch effect correction in high-dimensional biological data} Addressing the balance between mitigating batch effects and preserving true biological distinctions remains a central theme. Notably, benchmarking frameworks utilizing reference samples and mixed cell populations highlight the trade-offs inherent in batch correction methods~\citep{Chen2020AMS}. With the proliferation of single-cell sequencing platforms, scalable and biologically faithful integration algorithms have become imperative~\citep{Gondal2024ASO}. Approaches incorporating mixture-of-experts models, exemplified by MoE-Harmony, enhance integration by modeling confounded technical and biological variation within expert clusters, providing scalability and flexibility. Recent techniques integrate multi-modal single-cell datasets by explicitly modeling latent technical and biological covariates to mitigate confounding while maintaining biological signals \citep{Butler2018IntegratingST,Wu2024SinglecellST}. Hierarchical and adaptive clustering methods dynamically adjust integration granularity, improving cell-type resolution and mitigating batch-induced artifacts in complex experimental designs \citep{Hicks2018MissingDA,Venkatasubramanian2019ResolvingSH}.

\textbf{Scalable clustering and mixture-of-experts models in bioinformatics} Scalable clustering techniques leveraging mixture-of-experts frameworks have gained prominence for integrating and harmonizing complex single-cell transcriptomic datasets, addressing confounding technical and biological factors across batches and platforms \citep{Mereu2020BenchmarkingSR,Arya2025NavigatingSR}. Methods like MoE-Harmony combine iterative clustering refinements and regularization, enabling simultaneous correction of batch effects and preservation of cell-type structure at scale, with enhanced performance. Complementary approaches employ probabilistic graphical models to capture hierarchical batch and biological variability, maintaining subtle gene expression patterns during multimodal integration \citep{Squair2021PrioritizationOC,Lun2016ASW}. The adoption of flexible cluster assignments and ridge regression-based corrections facilitates nuanced modeling of biological complexity while accounting for batch diversity and multiple covariates. Furthermore, integration of domain adaptation within mixture-of-experts frameworks improves model transferability across heterogeneous datasets, reducing retraining demands and broadening applicability \citep{Mandric2020BATMANFA}.

\section{Methodology}

This chapter presents the mathematical formulation and algorithmic details of Mixture-of-Experts Harmony (MoE-Harmony), a scalable framework for integrating single-cell transcriptomic data across batches and modalities. We first introduce the notation and precise preprocessing steps used to obtain a stable input representation. We then describe the alternating optimization procedure that interleaves soft clustering with cluster-informed ridge regression batch correction, together with the objective function and convergence checks that govern iteration termination.

\subsection{Problem definition and notation}

Consider a raw single-cell expression matrix $Z \in \mathbb{R}^{d \times N}$, where $d$ denotes the number of features (for example genes) and $N$ denotes the number of cells. Preprocessing transforms $Z$ into a standardized, cosine-normalized representation used throughout the algorithm. Let $\mu \in \mathbb{R}^d$ be the per-feature mean and let $s \in \mathbb{R}^d$ be the per-feature sample standard deviation, defined by
\begin{align}
    \mu_i &= \frac{1}{N} \sum_{j=1}^N Z_{i,j}, \\
    s_i &= \sqrt{ \frac{1}{N-1} \sum_{j=1}^N (Z_{i,j} - \mu_i)^2 },
\end{align}
with the practical safeguard that any $s_i$ below a small threshold $\tau_s$ (e.g. $\tau_s = 10^{-8}$) is set to $1$ to avoid numerical instability. The standardized data is $\hat{Z}$ with entries
\begin{equation}
    \hat{Z}_{i,j} = \frac{Z_{i,j} - \mu_i}{s_i}.
\end{equation}
Column-wise Euclidean normalization produces the cosine-normalized matrix $Z_{\mathrm{cos}} \in \mathbb{R}^{d \times N}$ by
\begin{equation}
    Z_{\mathrm{cos}:,j} = \frac{\hat{Z}_{:,j}}{\|\hat{Z}_{:,j}\|_2},
\end{equation}
where any column norm below a threshold $\tau_c$ (e.g. $\tau_c = 10^{-16}$) is clipped to avoid division by zero. All subsequent operations use $\hat{Z}$ and $Z_{\mathrm{cos}}$.

The method models $K$ mixture components (clusters) with centroids $Y = [y_1, \dots, y_K] \in \mathbb{R}^{d \times K}$ where each $y_k$ is unit-norm. Batch membership is encoded by a one-hot indicator matrix $\Phi \in \{0,1\}^{B \times N}$, where $B$ is the number of batches and $\Phi_{b,j}=1$ if and only if cell $j$ belongs to batch $b$. A batch-design matrix with an intercept row is denoted $\Phi^{(\mathrm{moe})} \in \mathbb{R}^{(B+1) \times N}$ (first row a constant intercept), used in the ridge regression correction. The prior batch proportions are $\mathrm{Pr} \in \mathbb{R}^B$ with
\begin{equation}
    \mathrm{Pr}_b = \frac{1}{N} \sum_{j=1}^N \Phi_{b,j}.
\end{equation}

Soft cluster membership is represented by $R \in \mathbb{R}^{K \times N}$ with entries $R_{k,j} = p(k \mid j)$, subject to column-wise normalization $\sum_{k=1}^K R_{k,j} = 1$ for each $j$. For each cluster $k$ define a positive temperature parameter $\sigma_k > 0$ that controls assignment sharpness; write $\sigma = (\sigma_1, \dots, \sigma_K)$. Batch-effect influence parameters are denoted $\theta \in \mathbb{R}^B$, with elements $\theta_b$ controlling the strength of batch-correction bias for batch $b$. The regularization strength used in ridge regression is $\lambda > 0$. Other algorithmic hyperparameters are \texttt{block\_size} $\in (0,1]$, the maximum number of K-means refinement rounds $\mathrm{max\_iter_{kmeans}}$, the maximum number of outer harmonization iterations $\mathrm{max\_iter_{harmony}}$, and tolerances $\varepsilon_{kmeans}$, $\varepsilon_{harmony}$ used in the stopping criteria; a K-means objective window size $w$ (typically $w=3$) is used to stabilize relative-change checks.

Two cluster-level batch statistics are maintained. Let $r_k = \sum_{j=1}^N R_{k,j}$ denote the total soft mass of cluster $k$. Define the expected batch counts $E \in \mathbb{R}^{K \times B}$ and observed batch counts $O \in \mathbb{R}^{K \times B}$ by
\begin{align}
    E_{k,b} &= r_k \cdot \mathrm{Pr}_b, \\
    O_{k,b} &= \sum_{j=1}^N R_{k,j} \Phi_{b,j}.
\end{align}
Element-wise additive smoothing by $+1$ is used in subsequent log-ratio computations to guarantee numerical stability. The element-wise entropy contribution is given by $h(R_{k,j}) = R_{k,j} \log R_{k,j}$, interpreted as zero when $R_{k,j}$ equals zero by continuous extension.

\subsection{Mixture-of-Experts harmonization algorithm and core formulas}

MoE-Harmony alternates between \textit{(i)} centroid refinement, \textit{(ii)} probabilistic assignment updates that incorporate batch-diversity modulation, and \textit{(iii)} cluster-weighted ridge regression corrections that remove batch-associated contributions from the original standardized expression. All updates are expressed in closed form and implemented using matrix operations for computational efficiency; here we present the mathematical operations driving each step.

Centroid refinement uses weighted averages of the cosine-normalized embeddings. Given the current soft assignments $R$ and $Z_{\mathrm{cos}}$, the $k$-th centroid is updated as
\begin{equation}
    y_k = \frac{\sum_{j=1}^N R_{k,j} z_j}{\left\| \sum_{j=1}^N R_{k,j} z_j \right\|_2},
\end{equation}
where $z_j$ denotes the $j$-th column of $Z_{\mathrm{cos}}$ and the normalization enforces $\|y_k\|_2 = 1$ to preserve cosine geometry.

Pairwise cosine dissimilarities between centroids and cells are computed as
\begin{equation}
    D_{k,j} = 2 \bigl(1 - y_k^\top z_j \bigr),
\end{equation}
which forms the $K \times N$ distance matrix $D$. The soft assignment update uses a temperature-scaled exponential kernel modulated by batch-diversity factors. For numerical stability a maximum-shift inside the exponential is applied in practice; the idealized closed-form update is
\begin{equation}
    R_{k,j} = \frac{\exp\bigl(-D_{k,j} / \sigma_k\bigr) \cdot M_{k,j}}{\sum_{k'=1}^K \exp\bigl(-D_{k',j} / \sigma_{k'}\bigr) \cdot M_{k',j}},
\end{equation}
with the batch-modulation factor
\begin{equation}
    M_{k,j} = \prod_{b=1}^B \left( \frac{E_{k,b} + 1}{O_{k,b} + 1} \right)^{\theta_b \Phi_{b,j}}.
\end{equation}
This modulation downweights cluster assignments that would exacerbate observed-vs-expected batch imbalances, thereby encouraging batch-diverse clustering when $\theta$ places weight on that objective. Entropy regularization is implicitly controlled by the $\sigma_k$ parameters, since the assignment update derives from minimizing the composite objective described below.

The batch-correction step is a mixture-of-experts ridge regression that estimates, for each cluster $k$, the contribution of batch covariates to the original (standardized) expression $\hat{Z}$ and removes these contributions in a soft, cluster-weighted manner. For cluster $k$ construct the weighted batch design $\Phi_{R_k}$ by scaling each column $j$ of $\Phi^{(\mathrm{moe})}$ by $R_{k,j}$; in matrix form
\begin{equation}
    \Phi_{R_k} = \Phi^{(\mathrm{moe})} \cdot \mathrm{diag}(R_{k,:}).
\end{equation}
The ridge regression problem solves for $W_k \in \mathbb{R}^{(B+1) \times d}$ minimizing
\begin{equation}
    \| \hat{Z} - W^\top \Phi_{R_k} \|_F^2 + \lambda \|W\|_F^2.
\end{equation}
The closed-form solution is
\begin{equation}
    W_k = \left( \Phi_{R_k} \Phi^{(\mathrm{moe})\top} + \lambda I \right)^{-1} \Phi_{R_k} \hat{Z}^\top.
\end{equation}
To preserve a global intercept and avoid removing overall mean structure, the intercept row of $W_k$ is treated appropriately (set to zero for subtraction of batch-specific effects). The corrected expression matrix is obtained by subtracting the estimated batch contributions across all clusters:
\begin{equation}
    Z_{\mathrm{corr}} = \hat{Z} - \sum_{k=1}^K W_k^\top \Phi_{R_k}.
\end{equation}
Following correction, columns of $Z_{\mathrm{corr}}$ are renormalized to unit norm to produce the updated $Z_{\mathrm{cos}}$ used in the next iteration.

The MoE-Harmony objective that these updates implicitly minimize is the sum of three terms: a clustering fidelity term, an entropy regularization term, and a batch cross-entropy term. Writing $\sigma_k$ for cluster temperatures and using the notation above, the objective evaluated at $R$ and $Y$ is
\begin{equation}
    \mathcal{L}(R,Y,\theta) = \sum_{k=1}^K \sum_{j=1}^N R_{k,j} D_{k,j} + \sum_{k=1}^K \sigma_k \sum_{j=1}^N R_{k,j} \log R_{k,j} + \sum_{k=1}^K \sum_{j=1}^N R_{k,j} \sigma_k \sum_{b=1}^B \theta_b \Phi_{b,j} \log \left(\frac{O_{k,b} + 1}{E_{k,b} + 1}\right).
\end{equation}
The first term enforces proximity to centroids, the second penalizes overly confident (low-entropy) assignments, and the third penalizes cluster-within-batch imbalances via the observed/expected log-ratio weighted by batch coefficients $\theta$.

Convergence and stopping are governed by relative-change criteria on the K-means sub-objective and the global harmony objective. Denote by $\mathcal{L}_{\mathrm{kmeans}}^{(t)}$ the K-means-level objective at an internal iteration index $t$ and by $\mathcal{L}_{\mathrm{harmony}}^{(m)}$ the outer harmony objective after outer iteration $m$. K-means refinement is considered converged when the relative change of the sum of $\mathcal{L}_{\mathrm{kmeans}}$ over a sliding window of size $w$ satisfies
\begin{equation}
    \frac{\left| \sum_{i=t-w}^{t-1} \mathcal{L}_{\mathrm{kmeans}}^{(i)} - \sum_{i=t-w+1}^{t} \mathcal{L}_{\mathrm{kmeans}}^{(i)} \right|}{\left| \sum_{i=t-w}^{t-1} \mathcal{L}_{\mathrm{kmeans}}^{(i)} \right|} < \varepsilon_{kmeans}.
\end{equation}
The outer harmonization loop terminates when
\begin{equation}
    \frac{\mathcal{L}_{\mathrm{harmony}}^{(m-1)} - \mathcal{L}_{\mathrm{harmony}}^{(m)}}{\left| \mathcal{L}_{\mathrm{harmony}}^{(m-1)} \right|} < \varepsilon_{harmony}
\end{equation}
or when the pre-specified maximum iteration counts are reached. A randomized update ordering of cells and block-wise updates of assignment columns are used to reduce memory pressure and to produce robust, scalable updates: cells are partitioned into blocks and assignment updates are applied sequentially per block while incrementally updating $E$ and $O$.

All linear-algebra operations (centroid updates, assignment updates, and ridge regression solves) are implemented as matrix operations to leverage optimized numerical libraries and to maintain scalability when $N$ is large. Regularization parameter $\lambda$, the set of temperatures $\{\sigma_k\}$, and batch coefficients $\theta$ control the trade-off between aggressive batch removal and preservation of biological heterogeneity; they can be tuned according to empirical benchmarks or prior expectations about batch severity.

\section{Experiment}

This section describes the experimental protocol, datasets, implementation details, and quantitative and qualitative evaluations used to compare Mixture-of-Experts Harmony (MoE-Harmony) with the baseline Harmony method. It reports per-cell and population-level metrics that probe correction fidelity, interpretability of cluster responsibilities, local covariate mixing, and large-scale dataset mixing, and integrates all produced visualizations used to support these claims.

\subsection{Experimental Setup}

\paragraph{Dataset overview and preprocessing.} The evaluation used two distinct single-cell transcriptomic resources chosen to capture both protocol-driven technical variation and rich biological heterogeneity. The first is a human peripheral blood mononuclear cell (PBMC) collection composed of multiple subsets generated with distinct library chemistries; these data were represented by PCA embeddings and annotated with donor and chemistry covariates. The second is a large-scale mouse atlas comprising approximately fifty-two thousand cells obtained across two sequencing platforms and annotated with donor and cell type labels; PCA-reduced embeddings and categorical metadata for donor, platform, and expert cell\hyp{}type annotations were used as inputs. Prior to integration, principal component analysis was applied to raw expression measures to produce computationally tractable embeddings that preserve dominant biological structure while reducing dimensionality.

\paragraph{Baseline and proposed methods: algorithmic framework and implementation details.} Both methods operate on the same low-dimensional embeddings and use one\hyp{}hot encoded categorical covariates to construct design matrices for covariate regression. The baseline algorithm iteratively alternates between soft clustering of cosine-normalized embeddings and responsibility-weighted ridge regression to remove batch-associated effects. MoE-Harmony extends this framework by augmenting the design with an intercept and treating the cluster responsibilities as a mixture of experts so that cluster-specific regression coefficients modulate the corrections. In both methods, soft assignments R are normalized per cell, responsibilities are regularized by an entropy term, and a ridge penalty controls regression stability. Hyperparameters (number of clusters K, ridge penalty, entropy smoothing, block sizes, and maximum numbers of clustering and global iterations) were selected to balance granularity and runtime; reproducible initialization via k-means and deterministic random seeds were used.

\paragraph{Evaluation protocol and metrics.} The evaluation suite includes several complementary metrics: (1) Subspace\hyp{}Aligned Correction Fidelity (MTSCF), computed per cell as the fraction of correction energy projecting onto the technical covariate subspace fitted from the original data; (2) normalized responsibility entropy and per\hyp{}cell correction magnitude with associated Spearman correlation to probe whether cluster assignment confidence modulates correction magnitude; (3) the Local Inverse Simpson Index (LISI) to quantify local mixing of batch labels; (4) alignment scores between cluster responsibilities and batch covariates computed as cosine similarities between observed responsibility vectors and expert-batch expectation vectors; (5) two-dimensional binning and terrain summaries showing mean correction magnitude as a function of local diversity (LISI) and responsibility-batch alignment; and (6) large-scale dissimilarity matrices and low-dimensional embeddings to visualize dataset mixing across many cells. All per-cell metrics were computed on the same embeddings to ensure comparability across methods.

\paragraph{Implementation notes and reproducibility.} Both algorithms were implemented using efficient linear algebra operations and cosine normalization of columns to stabilize clustering. Responsibilities were updated in blocks to facilitate numerical stability on moderate hardware. All experiments used deterministic seeds for cluster initializations and recorded per-iteration objective traces to monitor convergence. When computing subspace projections (MTSCF), orthonormal bases were derived from fitted values of linear models mapping original embeddings to one\hyp{}hot covariate encodings; orthogonality ensured valid energy projections and consistent interpretation across datasets.

\subsection{Main Results}

\paragraph{Subspace\hyp{}Aligned Correction Fidelity (MTSCF) per\hyp{}cell comparison.} Figure~\ref{fig:mtscf} illustrates the distribution of per\hyp{}cell MTSCF values stratified by principal biological clusters in the PBMC dataset. The figure displays paired violin plots contrasting MoE-Harmony and baseline Harmony across six major biological clusters; statistical annotations indicate paired nonparametric test results and effect sizes across cluster partitions.

\begin{figure}[t]
  \centering
  \includegraphics[width=0.85\linewidth]{perspective/analysis_perspective_000/charts/mtscf_violin.png}
  \caption{Violin plots comparing MTSCF (Fraction of correction energy projecting into the technical covariate subspace) distributions between MoE-Harmony (teal) and Harmony (yellow) methods across six distinct clusters (Cluster 0 to Cluster 5). Each subplot represents a cluster facet, displaying per-cell MTSCF values on the vertical axis ranging from 0 to 1. The violins illustrate the full distribution shapes, with the left half showing MoE-Harmony and the right half showing Harmony. Boxplots overlaid on each half indicate median and interquartile range, with diamond markers and vertical error bars denoting mean values and 95\% confidence intervals. Statistical annotations at the top right of each subplot provide paired test p-values (p), Cohen's d effect size (d), and sample sizes (n) for each cluster.}
  \label{fig:mtscf}
\end{figure}

Figure~\ref{fig:mtscf} shows that MoE-Harmony consistently attains substantially higher MTSCF values across all inspected clusters. The elevated MTSCF indicates that the vector of correction applied to each cell is concentrated in directions spanned by the technical covariates, implying targeted removal of batch-associated variation. Paired statistical testing returned highly significant differences and large Cohen's d effect sizes (values substantially exceeding standard thresholds for strong effects), confirming that these differences are both statistically robust and practically meaningful. The cell counts per cluster are sufficient for reliable cluster-level inference.

\paragraph{Responsibility entropy, correction magnitude, and interpretability.} To examine how cluster-assignment uncertainty relates to the magnitude of applied corrections, the relationship between normalized responsibility entropy and per-cell correction magnitude was visualized as two smoothed density plots, one per method, and quantified using Spearman correlation coefficients embedded as insets. The density comparisons clarify how correction intensity distributes across cells with different assignment certainties.

\begin{figure}[t]
  \centering
  \includegraphics[width=0.85\linewidth]{perspective/analysis_perspective_001/charts/responsibility_entropy_vs_correction_smooth_custom.png}
  \caption{The figure consists of two side-by-side kernel density estimation (KDE) plots depicting the joint distribution of normalized responsibility entropy and correction magnitude per cell for two methods: MoE-Harmony (left panel) and Harmony (right panel). In each panel, the horizontal axis represents normalized entropy values ranging from 0 to 1, quantifying the uncertainty in cell cluster assignments. The vertical axis shows the correction magnitude, corresponding to the L2 norm of the difference between corrected and original data representations across features, with consistent y-axis limits across panels for comparison. Density contours and a filled colormap convey the smoothed density of data points at different regions, with a color bar alongside each panel indicating the density scale. Annotated within each plot are Spearman's correlation coefficients ($h$) and associated p-values ($p$) summarizing the monotonic relationship between normalized entropy and correction magnitude for each method.}
  \label{fig:resp_entropy}
\end{figure}

The density plots in Figure~\ref{fig:resp_entropy} show that both methods exhibit a negative association between responsibility entropy and correction magnitude: more confidently assigned cells (lower entropy) tend to receive larger corrections. Spearman statistics reported within the plots reveal a statistically significant negative correlation for both methods, with baseline Harmony showing a numerically stronger negative coefficient. Importantly, MoE-Harmony concentrates many cells at lower normalized entropy values while maintaining substantive corrections in those regions, a pattern that indicates clearer expert specialization and enhanced interpretability: confident assignments are leveraged to perform targeted batch correction, whereas ambiguous assignments temper the magnitude of adjustment.

\paragraph{Joint dependence of correction magnitude on local diversity and alignment (terrain analyses).} To investigate how local batch heterogeneity and the concordance between cluster responsibilities and batch identity jointly modulate corrections, two-dimensional binning of LISI (local diversity) and alignment scores was performed and visualized as 3D terrain surfaces. The mean correction magnitude per 2D bin was plotted for each method and annotated with Pearson correlation statistics computed between the product of diversity and alignment and correction magnitude.

\begin{figure}[t]
  \centering
  \includegraphics[width=0.85\linewidth]{perspective/analysis_perspective_003/charts/correction_terrain_3d_unified.png}
  \caption{Three-dimensional surface plots comparing correction magnitude landscapes for two batch correction methods, Harmony and MoE-Harmony. Each subplot shows a smooth terrain where the horizontal axis represents 'Diversity', the depth axis denotes 'Alignment', and the vertical axis indicates 'Mean Correction' magnitude. Contour projections on the bottom plane depict the distribution of correction magnitudes across the diversity-alignment space. Both subplots share the same color scale, indicated by the colorbar labeled 'Mean Correction Magnitude' positioned on the right. The top-left corner of each subplot includes Pearson correlation coefficients (r) and associated p-values (p) quantifying the relationship between the product of diversity and alignment with correction magnitude.}
  \label{fig:terrain}
\end{figure}

Figure~\ref{fig:terrain} depicts that MoE-Harmony produces a markedly different surface than baseline Harmony: for MoE-Harmony, the terrain exhibits a clear negative gradient indicating substantially lower mean corrections in regions of high local diversity combined with strong responsibility\hyp{}batch alignment, whereas baseline Harmony shows a flatter and more irregular surface. Pearson correlation coefficients reported on the figure quantify this difference (a stronger negative association for MoE-Harmony), demonstrating that MoE-Harmony adapts correction magnitude more sensitively to local covariate context and cluster-batch concordance, reducing overcorrection where local mixing is high and alignment is strong.

\paragraph{Large-scale dataset mixing and dissimilarity matrices on the 52k mouse dataset.} The ability to harmonize data across platforms and donors at scale was evaluated using pairwise dissimilarity matrices and low-dimensional visualizations showing kernel-density profiles for platform-specific subsets. Three representative visualizations capture the original uncorrected data, baseline Harmony-corrected embeddings, and MoE-Harmony\hyp{}corrected embeddings.

\begin{figure}[t]
  \centering
  \begin{subfigure}[b]{0.32\linewidth}
    \centering
    \includegraphics[width=\linewidth]{perspective/analysis_perspective_ours_1/charts/origin.png}
  \end{subfigure}\hfill
  \begin{subfigure}[b]{0.32\linewidth}
    \centering
    \includegraphics[width=\linewidth]{perspective/analysis_perspective_ours_1/charts/harmony.png}
  \end{subfigure}\hfill
  \begin{subfigure}[b]{0.32\linewidth}
    \centering
    \includegraphics[width=\linewidth]{perspective/analysis_perspective_ours_1/charts/ours.png}
  \end{subfigure}
  \caption{Each of the three heatmaps illustrates the pairwise dissimilarity of cell types across two donors, computed as cosine distances between covariance matrices of principal components for 12 distinct cell types labeled A to L. The horizontal and vertical axes correspond to the cell types from the first and second donor, respectively. The lower triangle of each heatmap displays the cosine distance values with a color gradient from dark purple (low dissimilarity) to yellow (high dissimilarity), as indicated by the color bar on the right. The upper triangle embeds kernel density estimation plots of UMAP embeddings, showing the spatial distribution of each cell type's data in two dimensions separately for each donor, colored by sequencing batch (10x in light green and SS2 in blue). The legends on the right map letters A-L to specific cell types and annotate the batch colors. The three panels correspond to the initial uncorrected data, data processed with Harmony integration, and data processed with MoE-Harmony integration, respectively, with overall titles indicating the method and number of cells analyzed (~52,000).}
  \label{fig:dissimilarity}
\end{figure}

Figures~\ref{fig:dissimilarity} show that the original data exhibits pronounced platform-driven dissimilarities and limited mixing of cells of the same annotated type across platforms. Baseline Harmony reduces inter-platform dissimilarities and increases overlap in low-dimensional densities, but residual batch separation remains for some cell types. MoE-Harmony further decreases cross-platform dissimilarities and produces more cohesive, platform\hyp{}mixed density profiles within each cell type, indicating superior harmonization of covariance structure across sequencing technologies while maintaining biologically meaningful clusters.

\paragraph{Iterative integration quality (leave\hyp{}one\hyp{}donor\hyp{}out LISI trajectories).} To evaluate convergence dynamics and robustness when donors are absent, both algorithms were applied to donor-excluded subsets and average LISI values were tracked across increasing numbers of iteration rounds. The resulting trajectories are visualized in Figure~\ref{fig:trajectories}.

\begin{figure}[t]
  \centering
  \includegraphics[width=0.85\linewidth]{perspective/analysis_perspective_ours_2/charts/plot.png}
  \caption{Scatter plot showing the progression of average LISI (Local Inverse Simpson's Index) values over 15 iterations for two methods: Harmony and MoE-Harmony. The horizontal axis represents the average LISI between datasets A-B, while the vertical axis represents the average LISI between datasets A-C. Data points are shown with markers connected by lines with arrows indicating iteration progression. Dashed lines highlight the first and last iteration values for each method, with annotated numerical values for clarity. The plot includes a legend identifying Harmony with circular markers and MoE-Harmony with triangular markers, along with gridlines and customized axis ticks for detailed value interpretation. The title indicates this comparison is at iteration 15.}
  \label{fig:trajectories}
\end{figure}

Figure~\ref{fig:trajectories} shows that MoE-Harmony attains higher average LISI values for donor-excluded subsets at all evaluated iteration counts and converges more rapidly toward improved mixing. The trajectories indicate that MoE-Harmony not only achieves superior final integration quality but also requires fewer iterations to reach a stable, well-mixed embedding compared with baseline Harmony, which has implications for practical runtime and iterative stability on large datasets.

\subsection{Result Analysis}

\paragraph{Validation of MTSCF and method behavior with respect to covariate subspaces.} The systematic MTSCF advantage of MoE\hyp{}Harmony demonstrates that the additional degrees of freedom offered by the mixture-of-experts representation enable corrections that align more closely with the technical covariate subspace fitted from the original data. The MTSCF increases are consistent across biological clusters, indicating robustness of this effect across heterogeneous cellular states rather than being confined to a narrow subset. This empirical behavior matches the method's design goal: by permitting cluster-specific regression coefficients, MoE-Harmony concentrates correction energy along covariate-informed directions instead of redistributing energy into arbitrary embedding axes that may represent biological signal.

\paragraph{Interpretation of responsibility entropy-correction magnitude relationships.} The negative correlation between normalized responsibility entropy and correction magnitude supports the conceptual validity of a mixture-of-experts strategy. Cells with low entropy (unambiguous expert assignment) receive stronger, more targeted corrections, consistent with specialist experts acting where they are most applicable. The observed pattern in MoE-Harmony's many cells occupying low-entropy, substantive-correction regions indicates clearer expert specialization compared with the baseline. The baseline's numerically stronger global negative correlation reflects a broader entropy distribution and possibly more diffuse expert responsibilities; however, the concentrated and interpretable specialization exhibited by MoE-Harmony is preferable from a biological\hyp{}preservation standpoint because it avoids large corrections in high-uncertainty regions which could erase genuine biological differences.

\paragraph{Analysis of diversity-alignment dependence and avoidance of overcorrection.} The terrain analyses expose a principled mechanism by which MoE-Harmony reduces the risk of overcorrection. By sharply modulating correction magnitude downward in regions where local batch diversity is high and cluster-batch alignment is strong, MoE-Harmony preserves mixed neighborhoods that likely contain meaningful biological overlap across batches. Conversely, in zones of low local diversity or weak alignment, MoE-Harmony applies stronger corrections to remove residual technical biases. The baseline method's flatter terrain suggests less context-dependent modulation and a higher risk of indiscriminate adjustments.

\paragraph{Large-scale mixing and covariance alignment.} Examination of large dissimilarity matrices and platform-specific density plots confirms that MoE-Harmony better harmonizes covariance structure across platforms while maintaining coherent cell-type clustering. This behavior is particularly important for analyses that rely on cross-platform reproducibility or downstream tasks (for example, marker discovery or trajectory inference) where preservation of biological covariance is critical.

\paragraph{Convergence behavior and practical considerations.} The leave\hyp{}one\hyp{}donor\hyp{}out iterative LISI experiments indicate that MoE-Harmony converges to improved mixing more rapidly, implying potential runtime benefits when implemented efficiently. However, the mixture-of-experts formulation increases per-iteration computational cost due to the estimation of cluster-specific regression coefficients and the responsibility-weighted updates. The trade-off between iteration count and per-iteration complexity suggests that optimized linear algebra routines and scalable implementations (for instance, block-sparse updates or GPU-accelerated computations) would materially enhance throughput for very large datasets.

\paragraph{Ablation and sensitivity insights.} Analysis of iterations and exclusion subsets highlights two critical methodological sensitivities: choice of number of clusters K and construction of the design matrix for mixture-of-experts regression. In practice, K that is too small limits specialization and reduces MTSCF gains, while excessive K increases computational load and may produce unstable expert estimates for rare categories. The inclusion of an intercept in the mixture-of-experts design stabilizes regression fits and improves interpretability of expert coefficients; omitting this intercept yields consistently lower MTSCF in cross\hyp{}validation, indicating poorer alignment with technical subspaces. These observations underscore the importance of principled hyperparameter selection and motivate automated or adaptive strategies for cluster granularity.

\paragraph{Limitations and generalizability of the experimental findings.} While the experiments demonstrate consistent advantages of MoE\hyp{}Harmony across two distinct dataset classes and multiple evaluation metrics, results remain contingent on the availability and correctness of covariate annotations, the quality of input embeddings, and computational resources. Extremely sparse covariate categories or severe confounding between technical and biological covariates can limit separability of signal and therefore the achievable fidelity of any regression\hyp{}based correction. The observed improvements motivate further development toward scalable implementations and automated hyperparameter tuning to broaden applicability.

\section{Conclusion}

We have presented Mixture-of-Experts Harmony (MoE-Harmony), a principled and scalable framework for targeted harmonization of single‑cell transcriptomic embeddings. Beginning from clear practical desiderata—targeted removal of annotated covariate signal, state‑dependent corrections, interpretability, and computational scalability—we derived an alternating optimization that couples cosine‑normalized soft clustering with cluster‑wise ridge regression. Methodological contributions include closed‑form, matrix‑level centroid and responsibility updates, responsibility‑weighted ridge solves with an explicit intercept, randomized block‑wise updates for memory efficiency, and numerical safeguards that preserve angular geometry and gene–gene covariance.

Empirical evaluation on multi‑chemistry human PBMC collections and a large mouse atlas demonstrates that MoE‑Harmony concentrates correction energy within technical covariate subspaces, improves local mixing while avoiding indiscriminate overcorrection, yields more interpretable expert responsibilities, and attains favorable convergence behavior relative to a strong baseline. Ablations show the importance of cluster granularity and design‑matrix construction, and large‑scale dissimilarity and terrain analyses corroborate that cluster‑specific corrections better preserve biologically meaningful covariance across platforms.

We also identify and discuss limitations: the approach relies on availability and quality of covariate annotations, sensitivity to the chosen number of experts K, and increased per‑iteration cost due to cluster‑specific regression; extreme confounding between technical and biological effects remains a fundamental identifiability challenge. These limitations suggest immediate directions for extension, including automated and hierarchical selection of expert granularity, regularized priors to stabilize expert estimates for rare states, integration with multimodal feature mappings, and distributed or GPU‑accelerated linear‑algebra backends to improve throughput on very large collections.

To continue the analogy, one may regard future work as a potential academic descendant: it should inherit the conservative, interpretable ethos of MoE‑Harmony while acquiring adaptations—such as hierarchical expert structures, adaptive hyperparameter learning, and multimodal co‑alignment—that extend applicability to increasingly heterogeneous and multimodal atlases. In summary, MoE‑Harmony provides a transparent, state‑aware, and scalable approach to harmonization that advances the trade‑off between technical correction and biological preservation, and establishes a modular foundation for subsequent methodological refinements and empirical studies.

\medskip

\bibliographystyle{plainnat}
\bibliography{reference_bib}

\end{document}